\chapter{可靠性：让失败可控、可恢复}
\label{chap:reliability}

\begin{takeaway}[学习目标]
把 Agent 失败分层，设计超时、重试、检查点、补偿、降级和恢复；理解概率组件如何与确定性护栏组合，并能写出服务级目标与故障演练计划。
\end{takeaway}

\section{生产可靠性不是“多试几次”}

Agent 的一次任务可能包含十几次模型与工具调用。即使每一步成功率是 98\%，二十个相互依赖步骤全部成功的理想概率也只有约 67\%。真实系统还会遇到网络、限流、数据变化、权限和并发问题。可靠性必须来自系统设计。

\section{先建立故障分类}

\begin{table}[htbp]
\centering
\small
\begin{tabularx}{\textwidth}{>{\bfseries}p{0.18\textwidth} X X}
\toprule
层级 & 例子 & 首选处理 \\
\midrule
模型 & 协议错误、选错工具、幻觉 & 校验、有限重试、换模型、拒答 \\
工具 & 超时、限流、依赖不可用 & 退避、熔断、替代工具 \\
状态 & 重复执行、版本冲突、丢检查点 & 幂等、事务、事件日志 \\
上下文 & 约束丢失、证据过期 & 重新组装、来源时间校验 \\
策略 & 越权、预算超限 & 权限层拒绝、停止、人工确认 \\
业务 & 目标本身不可完成 & 部分结果、澄清、升级人工 \\
\bottomrule
\end{tabularx}
\end{table}

错误分类应产生机器可读码。只有知道错误属于瞬时、永久、冲突还是需人工，恢复策略才不会盲目。

\section{超时预算从外向内分配}

如果用户请求最多等待 60 秒，内部步骤的超时总和不能无限扩张。为整体运行设置 deadline，工具每次调用从剩余时间获取子预算。达到 deadline 后取消未完成任务，返回已有结果与状态。

取消必须向下传播：模型流式请求、HTTP 调用、子 Agent 和后台任务都应响应取消。否则用户看见“已停止”，系统仍可能继续产生费用或副作用。

\section{重试矩阵}

\begin{itemize}
  \item 网络断开、HTTP 429/503：有限次数指数退避，遵守 \code{Retry-After}。
  \item schema 不合法：把具体校验错误反馈给模型，通常重试 1--2 次。
  \item 参数业务非法：不要原样重试，需修正参数或重新规划。
  \item 权限拒绝：禁止重试绕过，向用户说明所需权限。
  \item 写操作结果未知：先用幂等键查询状态，不直接重复执行。
\end{itemize}

全局重试预算能防止嵌套组件各自重试三次，最终放大成几十次调用。每次重试都记录原因、等待时间和新参数。

\section{检查点与恢复}

长任务在关键步骤后保存检查点：计划版本、已完成步骤、artifact 句柄、工具调用状态、剩余预算和待确认动作。恢复时从最后一个\term{已提交且可验证}的检查点继续，而不是重新播放所有自然语言消息。

检查点需要版本 schema；代码升级后应能迁移旧状态，或明确拒绝并提供安全重启。包含凭证的状态不要明文持久化，只保存短期引用。

\section{补偿事务}

跨多个外部系统的流程通常无法使用数据库事务。可以采用 Saga 思路：每个正向动作对应补偿动作。例如创建日历事件后发送通知，通知失败时可以删除事件或标记为待处理。

补偿不等于完美撤销。邮件发出后无法收回，价格变化也可能使退款产生差额。设计时要标明“可逆、可补偿、不可逆”三种语义，并把不可逆步骤尽量放在最后和确认之后。

\section{熔断、隔离与降级}

依赖持续失败时，熔断器暂时停止调用，避免拖垮整个系统。不同租户、工具和任务类型应有独立并发池，防止一个大任务占满资源。降级路径可以是：换备用检索源、使用缓存、转为只读、生成草稿而不提交，或交给人工。

降级必须对用户可见。系统不应在没有实时数据时假装结果是最新的，应标注数据时间和缺失能力。

\section{确定性外壳包住概率核心}

模型擅长理解与候选生成；代码擅长约束与检查。以下部分优先使用确定性实现：权限、金额计算、schema、状态迁移、去重、预算、停止条件、引用存在性和最终提交。模型可以提出计划，但不能修改这些规则。

\section{服务目标与错误预算}

不要只监控 API 是否返回 200。Agent 的服务目标可以包括：任务成功率、危险动作误执行率、P95 完成时间、单任务成本、恢复成功率和人工接管率。对高风险动作，安全指标优先于平均完成率。

错误预算帮助团队决定是否继续加功能。当回归失败或越权事件超过预算时，应暂停扩张，先处理工具与评测缺口。

\section{故障演练}

定期注入：模型超时、工具返回畸形 JSON、检索源为空、网络断开、重复 webhook、状态库短暂不可用、批准后资源已变化、子 Agent 永不返回。验证系统是否停止在可理解状态，并能从检查点恢复。

\begin{practice}[为最小 Agent 加恢复]
让每三次工具调用随机失败一次。加入错误分类、全局重试预算、幂等键、检查点和取消。运行 100 个任务，统计成功率、重复副作用次数、P95 步数和无法恢复的故障。
\end{practice}

\begin{checkpoint}
\begin{itemize}
  \item \checkbox 每种错误都有可恢复性分类，而非统一重试。
  \item \checkbox 写操作能用幂等键确认是否已经完成。
  \item \checkbox 长任务能从已验证检查点继续。
  \item \checkbox 降级和数据过期会明确展示给用户。
\end{itemize}
\end{checkpoint}

